\chapter{总结与讨论}
\label{chap:conclusion}

\section{总结}

本研究利用希腊及地中海地区的 FireCube 多变量时空数据集，进行野火风险预测的研究。首先介绍FireCube数据集的区域概况及时空数据构建，采样方式及训练数据，并详细分析其特征变量，为模型的构建奠定基础。

模型采用时空解耦的双分支架构：利用 3D Swin Transformer 分支通过三维滑动窗口机制提取气象序列的时空演化特征；利用 2D Swin Transformer 分支编码地形、植被等静态地理背景先验。在特征融合阶段，设计了地理感知交叉注意力机制（Geo-CA）。为弥补数据驱动模型在物理感知上的局限，本研究创新性地采用物理启发式特征重构。通过克劳修斯-克拉珀龙方程推导并利用马格努斯-泰登经验公式，将原始气温与露点温度变量重构为饱和水汽压差（VPD），用以定量刻画大气环境的潜在蒸发需求与植被燃料的脱水压力。此外，模型通过对 U（经向）和 V（纬向）风速分量的物理还原，得到风场矢量，并通过计算局部风场矢量与空间相对向量的点积，将风场动力学偏置硬编码至注意力权重计算中。这一设计打破了标准注意力机制的各向同性局限，使模型能够感知风场信息。

实验结果表明，该模型在测试集上的 F1-Score 达到 84.52\%，显著优于 LSTM、ConvLSTM 等主流基准模型。消融实验证实，饱和水汽压差VPD的特征重构有效降低了误报率，而风场偏置注入则显著提升了对下风向危险信号的捕捉能力（召回率达 84.58\%）。最后，通过 SHAP 归因分析验证了最大风速与 VPD 的对野火概率预测的突出贡献地位，并结合空间交叉注意力图的可视化，直观展示了注意力分布随物理风场偏移的规律，证实了模型决策逻辑与林火空气动力学常识的高度一致性。

\section{局限性及未来展望}

本研究虽然在融合物理信息的野火预测方面取得了一定进展，但仍存在若干局限性。首先，在数据表征维度上，目前输入的特征变量尚不充足，缺乏如火势历史动态轨迹等关键物理过程特征，且土地覆盖分类的粒度仍有提升空间，难以捕捉微观尺度的燃料结构差异。受限于原始数据集的分辨率，本研究采用的1km$\times$1km空间尺度在处理地形剧烈波动或城乡交界处的精细化蔓延时仍显粗糙。其次，本实验主要聚焦于地中海气候下的希腊地区，由于全球不同地区的火灾驱动机制存在区别，模型在跨气候带的大规模泛化能力及在全球范围内的训练样本储备仍需进一步加强。此外，受计算资源及样本规模限制，当前双分支网络的深度配置较为克制，在提取深层高阶物理演化语义方面可能存在表征瓶颈。

展望未来，研究工作将重点从以下几个方向展开。一是多源高分辨率数据的深度集成，通过引入遥感影像数据及具备穿透能力的雷达数据，提升模型对复杂环境的感知精度。二是探索全球化的预训练与迁移学习范式，利用海量多区域数据进行自监督学习，以克服单一区域正样本不足的局限性。三是进一步深化物理信息与深度学习架构的耦合，尝试将更为复杂的野火蔓延方程以算子学习或损失函数约束的形式嵌入更深层的模型架构中，旨在构建参数容量更大、物理一致性更强且具备全球预测能力的野火风险预测大模型。